% =====================================================================
%  Capítulo 6 — Dados, modelo e inferência embarcada (só na AP3)
% =====================================================================
\chapter{Pipeline de dados, modelo e inferência embarcada}
\label{cap:modelo}
\rubricaFuncionalidade

\begin{orientacao}
Este capítulo só é compilado quando \texttt{\textbackslash def\textbackslash CICLO\{3\}}.
Ele documenta o projeto integrador de AIoT: do dado bruto ao modelo
rodando no dispositivo restrito.
\end{orientacao}

\section{Pipeline de dados}

\begin{orientacao}
Sincronização, rotulagem, janelamento e extração de características.
Diga o tamanho da janela e a sobreposição, e justifique ambos pelo
fenômeno que está sendo classificado --- não pelo que era o padrão da
ferramenta.
\end{orientacao}

\section{Modelo e treinamento}

\begin{orientacao}
Arquitetura, divisão treino/validação/teste (com a regra de separação:
por tempo? por indivíduo? aleatória? a escolha muda tudo), métricas e
matriz de confusão. Acurácia sozinha não descreve um classificador
desbalanceado.
\end{orientacao}

\begin{table}[H]
\centering
\caption{Desempenho do modelo por classe}
\label{tab:metricas}
\footnotesize
\begin{tabular}{@{}L{4.0cm}R{2.6cm}R{2.6cm}R{2.6cm}R{2.0cm}@{}}
\toprule
\textbf{Classe} & \textbf{Precisão} & \textbf{Revocação} & \textbf{F1} & \textbf{Amostras} \\
\midrule
\ldots & 0,00 & 0,00 & 0,00 & 0 \\
\bottomrule
\end{tabular}
\fonte{Elaborada pela equipe.}
\end{table}

\section{Otimização para o dispositivo restrito}

\begin{orientacao}
Quantização, \emph{pruning} ou destilação: o que foi aplicado, e quanto
custou em acurácia. A tabela abaixo é o coração da \rotuloAP\ ---
ela mostra o preço de cada grama de memória economizada
\cite{warden2019}.
\end{orientacao}

\begin{table}[H]
\centering
\caption{Custo da otimização: antes e depois}
\label{tab:otimizacao}
\footnotesize
\begin{tabular}{@{}L{4.2cm}R{2.8cm}R{2.8cm}R{3.2cm}@{}}
\toprule
\textbf{Métrica} & \textbf{Modelo original} & \textbf{Modelo embarcado} & \textbf{Variação} \\
\midrule
Acurácia & 0,00 & 0,00 & --- \\
Tamanho em \emph{flash} (kB) & 0 & 0 & --- \\
RAM de pico (kB) & 0 & 0 & --- \\
Latência de inferência (ms) & --- & 0 & --- \\
Energia por inferência (mJ) & --- & 0 & --- \\
\bottomrule
\end{tabular}
\fonte{Elaborada pela equipe.}
\end{table}

\section{Inferência no dispositivo}

\begin{orientacao}
Como o modelo foi embarcado, quanto ele ocupa, quanto tempo leva uma
inferência medida no próprio ESP32 e o que o sistema faz com o
resultado. Fechem o ciclo: a decisão de atuação da \rotuloAP\ 1 agora é
tomada pelo modelo?
\end{orientacao}
